% English abstract. Faithful translation of the German Kurzfassung
% (kurzfassung.tex), which follows on the next page.

\section*{Abstract}

\begingroup
\small
\setlength{\parskip}{0.35em}

Large language models can justify their answers in natural language.
However, the faithfulness of such explanations is not guaranteed. They can
sound plausible and still conceal the factors that actually determined the
answer. The Causal Concept Faithfulness Metric by Matton
et al.~\cite{mattonWalkTalkMeasuring2024} measures this faithfulness by
counterfactually altering one concept of the input at a time and comparing the
measured causal effect with the effect the explanation implies. Because only
one concept is altered at a time, the metric can underestimate or misattribute
the influence of interacting concepts. This happens, for example, when a
removed concept can be reconstructed from another one or when only the
combination of several concepts changes the answer.

As a solution, this thesis develops FELICS (Faithfulness Estimation Leveraging
Interacting Concept Sets). An oracle combines concepts for which an interaction
is plausible into bounded candidate groups. Within each group, every nonempty
combination of concepts is altered jointly. Shapley values distribute the
measured coalition effects across the individual concepts, so the comparison
with the explanation still takes place at the concept level. A reference-based
correction makes effects from single and joint interventions comparable, and
an upper bound on the group size limits the number of counterfactuals needed.
For an imposed bound of 1, FELICS formally reduces to the original metric.

FELICS was implemented and evaluated on the BBQ and MedQA datasets and on the
newly created tabular dataset New German Credit. The models studied were
GPT-OSS 120B and Mistral Medium 3.5, with 30 questions per dataset each. The
evaluation covers four research questions and a case study on the MedQA example
by Matton et al.~\cite{mattonWalkTalkMeasuring2024} that motivated this thesis.

The results show no general relationship between the strength of the
interaction suspected by the oracle and the change in causal effects. In
separately executed runs, the theoretical compatibility with the original
metric could not be confirmed empirically within the predefined bounds. The
correction largely preserved the ranking of concepts in BBQ and MedQA but
reordered it substantially in New German Credit. Joint interventions did not
systematically increase the measured faithfulness. The direction of the change
depended on the dataset and the model, and at the dataset level faithfulness
changed only slightly. In the case study, only one of the two concepts examined
received a higher causal effect, although both had been expected to. At the
same time, individual examples show that joint interventions can reveal
threshold-like behavior that singleton interventions conceal, and that they can
limit double counting when the same evidence is spread over overlapping
concepts.

The contribution of this thesis is to introduce joint interventions into the
measurement of faithfulness. In this way, FELICS makes visible how joint
interventions redistribute causal influence between concepts. It is intended
as a practical, interaction-aware diagnostic tool and not as a method that
determines a unique true interaction structure.

\endgroup
